\section{端到端链路：从模型文件到下一个 token}

\subsection{先把一条请求链路钉住}

推理引擎里有很多对象：manifest、typed graph、量化、KV cache、CPU/GPU backend、oracle、profiler。若不先把它们放进同一条请求链路，读者很容易学到一堆孤立概念。这里用一个最小对话请求把全链路钉住：用户输入一段 prompt，引擎加载模型，编译计划，执行 prefill，进入 decode 循环，最后不断产生 token。

\begin{lstlisting}[numbers=none,caption={端到端推理链路的最小形态}]
model directory
  -> manifest
  -> verified typed graph
  -> quantization and layout plan
  -> CPU/GPU executable plan

request prompt
  -> tokenizer
  -> prefill plan
  -> KV cache initialized
  -> logits
  -> sampler
  -> decode plan loop
  -> next token stream
\end{lstlisting}

先把这些词落到一个能检查的最小目录里。假设磁盘上有一个教学模型目录：

\begin{lstlisting}[numbers=none,caption={最小模型目录里每个文件的责任}]
mini-model/
  config.json          hidden size, layers, heads, vocab size
  tokenizer.json       text <-> token id
  tensors.bin          raw weight bytes
  tensors.index.json   tensor name, dtype, shape, offset, checksum
  manifest.json        canonical model contract
\end{lstlisting}

\code{manifest} 不是“又一个配置文件”。它把外部文件名、权重命名、shape、dtype、量化格式、tokenizer 版本和校验和整理成引擎内部可以相信的合同。没有 manifest，后端只能靠字符串猜 \code{model.layers.0.attn.q.weight} 到底是不是第 0 层 Q projection；一旦模型方言换名，错误会延迟到运行时。manifest 能证明“这些文件和权重按本引擎认识的规则组成同一个模型”，但它不能证明算子实现正确，也不能证明生成文本质量好。

\code{typed graph} 是把 manifest 里的张量和配置连成有类型的计算图。最小 decoder 层里，它会记录：hidden 输入是 \code{[1,H]}，Q/K/V projection 输出是什么 shape，attention 读取和写入哪些 KV cache state，MLP 输出仍回到 \code{[1,H]}。typed graph 能证明 shape、dtype、layout 和状态边在图层面说得通；它不能证明某个 AVX2 kernel 足够快。

\code{KV cache} 是请求状态，不是模型权重。prefill 处理 prompt 时，把每一层每个历史 token 的 K/V 写进去；decode 第一步之后，每生成一个新 token，都要读取旧 K/V 并追加当前位置的 K/V。若把 KV cache 当成普通临时 activation，memory planner 可能复用或覆盖它，短 prompt 也许不炸，第二个 token 以后 logits 就会漂。

\code{backend} 是执行计划落到机器的选择：CPU scalar、AVX2、AVX-512、GPU kernel 或库调用。它只负责按合同执行 segment，不应该在热路径重新猜 shape。 \code{oracle} 是裁判，用 reference 或分层检查比较结果；\code{profiler} 是账本，记录 prepare、prefill、decode、copy、sampler 和 fallback 的成本。oracle 证明“结果在误差合同内”，profiler 证明“时间花在哪里”；二者不能互相替代。

这条链路有两个方向。第一条是“模型到计划”：它尽量在请求前完成，把能检查的语义、shape、dtype、量化、layout 和 kernel 选择提前固定。第二条是“请求到 token”：它发生在 runtime 中，绑定 session、token buffer、KV cache、arena、workspace 和 sampler。把这两条线分开，是推理引擎不退化成图解释器的关键。

\begin{keyidea}
推理引擎的核心不是“写几个快 kernel”，而是把模型资产编成可执行计划，再把计划安全地绑定到一次会增长的请求状态上。优化只能发生在这条链路上的某个明确位置。
\end{keyidea}

\subsection{从一个坏的 \code{generate()} 看边界}

最容易写出的推理程序，是一个巨大的 \code{generate()} 函数：读模型、解析 tokenizer、分配 activation、判断量化格式、选择 kernel、采样 token 全挤在一个循环里。

\begin{lstlisting}[numbers=none,caption={不可维护的推理主循环形态}]
generate(prompt):
  model = read_files()
  tokens = tokenize(prompt)
  for step in 0..max_new_tokens:
    for layer in model.layers:
      if layer.weight_format == "q4":
        unpack_and_run_q4()
      else:
        run_f32()
      maybe_allocate_more_memory()
      maybe_move_to_gpu()
    next = sample(logits)
    tokens.push_back(next)
\end{lstlisting}

问题不在于代码短，而在于它把四类生命周期混在一起：模型文件和权重属于模型生命周期；packed weight 和 executable plan 属于编译生命周期；token buffer、KV cache 和 sampler 属于请求生命周期；单个 kernel 的 workspace 和 profiler event 属于算子生命周期。边界混乱以后，错误定位、内存复用和性能测量都会失真。

\begin{center}
\begin{tabular}{p{0.20\linewidth}p{0.34\linewidth}p{0.31\linewidth}}
\toprule
生命周期 & 应该保存的对象 & 不应该做的事 \\
\midrule
模型生命周期 & 文件索引、manifest、原始权重视图 & 持有某次请求的 token 状态 \\
编译生命周期 & typed graph、packed weight、backend plan & 每个 token 重新解析模型结构 \\
请求生命周期 & session、KV cache、sampler、trace & 修改模型 manifest 或后端 capability \\
算子生命周期 & 局部 workspace、临时 tile、计数器 & 泄漏到下一层或下一次请求 \\
\bottomrule
\end{tabular}
\end{center}

\begin{warningbox}
推理引擎里很多难查的问题，本质上不是数学错，而是生命周期错：静态对象在热路径重复构造，请求对象被全局共享，临时 buffer 被跨 step 使用，或者 profiler 把 prepare 成本算进 decode。
\end{warningbox}

\subsection{模型到计划：哪些事情应该提前做}

模型加载阶段不应该只读文件。它要把外部资产变成编译器可以相信的事实。

\begin{center}
\begin{tabular}{p{0.22\linewidth}p{0.33\linewidth}p{0.30\linewidth}}
\toprule
阶段 & 产物 & 拦住的问题 \\
\midrule
读取资产 & raw tensor index、config、tokenizer & 文件缺失、大小不一致、checksum 错 \\
manifest & canonical tensor roles & 模型方言、权重命名、tied weight \\
verifier & typed graph 与诊断 & shape/dtype/quant/KV 合同错误 \\
优化与 lowering & executable plan & 融合、layout、memory、backend 选择 \\
prepare & packed weight、device buffer & 热路径重复 packing 或上传 \\
\bottomrule
\end{tabular}
\end{center}

这张表的核心是“早拒绝”。如果 tokenizer 词表大小和 embedding 行数不一致，应该在前端报错；如果 int4 scale 数量不匹配，应该在量化 verifier 报错；如果 GPU backend 不支持某种低比特格式，应该在 plan 阶段显式 fallback 或拒绝。不要等到 decode 第 128 个 token 时才用错误 logits 暴露问题。

\subsection{请求到 token：哪些事情只能运行时知道}

运行时不是编译器的替代品。它处理的是请求期才知道的状态：prompt 内容、当前 position、最大生成长度、sampler 参数、KV cache 已写入长度、当前选择的 backend、是否打开 oracle/profiler。

\begin{lstlisting}[numbers=none,caption={一次请求的运行期状态}]
Session:
  token_buffer
  current_position
  kv_cache
  activation_arena
  workspace
  sampler_state
  profiler_trace
  oracle_policy
\end{lstlisting}

prefill 和 decode 的职责不同。prefill 读取整个 prompt，批量写入 KV cache，并产出最后一个 prompt token 的 logits。decode 每次只处理一个新 token，读取历史 KV cache，追加当前位置的 K/V，再产出下一步 logits。很多后端优化之所以失败，是因为把这两种形态混成一个计划。

同一个算子名字在不同阶段不一定是同一种机器工作。QKV projection 在 prefill 里更像大 GEMM，在 decode 里更像小 batch GEMV；attention 在 decode 里还要读取从 0 到当前 position 的历史 K/V；sampler 若每步强制 GPU 同步并拷回完整 logits，也会进入端到端延迟。

\begin{systemdiagram}{模型编译链路与请求执行链路}
\begin{pdfdiagram}[x=1cm,y=1cm]
  \node[book row label] at (-0.25,2.65) {编译链};
  \node[book accent node,text width=2.0cm] (asset) at (1.0,2.65) {模型资产};
  \node[book teal node,text width=2.0cm] (manifest) at (3.55,2.65) {manifest};
  \node[book purple node,text width=2.0cm] (graph) at (6.1,2.65) {typed graph};
  \node[book amber node,text width=2.0cm] (plan) at (8.65,2.65) {executable\\plan};

  \node[book row label] at (-0.25,0.85) {请求链};
  \node[book accent node,text width=2.0cm] (prompt) at (1.0,0.85) {prompt\\tokens};
  \node[book teal node,text width=2.0cm] (prefill) at (3.55,0.85) {prefill};
  \node[book purple node,text width=2.0cm] (kv) at (6.1,0.85) {KV cache};
  \node[book amber node,text width=2.0cm] (decode) at (8.65,0.85) {decode\\loop};
  \node[book navy node,text width=2.0cm] (token) at (11.2,0.85) {next token};

  \draw[book strong arrow] (asset) -- (manifest);
  \draw[book strong arrow] (manifest) -- (graph);
  \draw[book strong arrow] (graph) -- (plan);
  \draw[book dashed arrow] (plan) -- (prefill);
  \draw[book dashed arrow] (plan) -- (decode);
  \draw[book strong arrow] (prompt) -- (prefill);
  \draw[book strong arrow] (prefill) -- (kv);
  \draw[book strong arrow] (kv) -- (decode);
  \draw[book strong arrow] (decode) -- (token);
\end{pdfdiagram}
\diagramnote{本图要证明的命题是：模型编译链和请求执行链共享计划，但生命周期和错误边界不同。}
\end{systemdiagram}

\subsection{计划不是抽象名词，而是一份可执行合同}

本册反复使用 \term{可执行计划}{executable plan} 这个词。它不是架构口号，而是一份后端和运行时都能执行的合同：语义 segment、kernel 选择、layout、packed weight、workspace、KV cache 读写、profiler event 和 fallback 规则都要落在这里。

\begin{lstlisting}[numbers=none,caption={executable plan 的最小合同字段}]
ExecutablePlan:
  segments:
    id
    operator_or_fused_region
    input_values
    output_values
    state_effects: reads_kv | writes_kv | none
    backend_choice
    kernel_descriptor
    workspace_bytes
    profiler_label
    oracle_checkpoint_policy
\end{lstlisting}

这里的 \code{state_effects} 很重要。decode attention 会读历史 KV cache，同时写入当前位置的 K/V；sampler 会读 logits 并更新随机数状态。若 IR 和 plan 不表达这些状态边，memory planner 可能错误复用 buffer，profiler 可能漏掉 cache 写入，oracle 也无法定位漂移点。

\subsection{一个真实业务接口长什么样}

说“能做业务开发”，不能只说能生成文本。最小可用接口要把请求参数、资源预算、流式输出、停止原因和诊断 trace 都暴露出来。一个本地推理服务的请求可以先设计成：

\begin{lstlisting}[numbers=none,caption={GenerateRequest 的最小字段}]
GenerateRequest:
  request_id
  model_id
  prompt
  max_new_tokens
  max_context_tokens
  temperature
  top_k
  top_p
  repetition_penalty
  stream: true | false
  quant_profile_id
  backend_policy: cpu | gpu | auto
  trace_level: none | summary | debug
\end{lstlisting}

这些字段不是接口装饰。 \code{max_context_tokens} 决定 KV cache 容量；\code{max_new_tokens} 决定最坏 decode 时长；采样参数决定结果可复现性；\code{quant_profile_id} 决定误差和性能合同；\code{backend_policy} 决定 CPU/GPU plan 选择；\code{trace_level} 决定是否允许 materialize 中间结果和更细 profiler。

响应也不能只是一个字符串。流式输出至少要区分 token 事件、错误事件和结束事件：

\begin{lstlisting}[numbers=none,caption={流式响应事件}]
TokenEvent:
  request_id
  token_id
  token_text
  index
  model_position
  finish_reason_optional
  trace_event_id_optional

FinalResponse:
  request_id
  text
  generated_tokens
  finish_reason: stop_token | max_tokens | context_limit | cancelled | error
  usage:
    prompt_tokens
    completion_tokens
    kv_cache_bytes
  trace_summary_optional
\end{lstlisting}

这样设计后，业务层可以实现超时、取消、计费、审计、问题定位和重试。若接口只返回一个最终字符串，就无法判断卡在 prefill、decode、sampler、KV 容量还是后端 fallback。

\subsection{流式生成不是 \code{print} 循环}

流式生成的状态机要明确。prefill 完成之前，通常还没有第一个输出 token；prefill 之后每个 decode step 产生一个候选 token，sampler 决定是否结束，然后服务层把 token 事件推给客户端。

\begin{lstlisting}[numbers=none,caption={流式生成状态机}]
state = admitted
tokens = tokenizer.encode(prompt)

run prefill(tokens)
emit trace: first_token_ready

while not stopped:
  logits = decode_one_token(session)
  next = sampler.sample(logits)
  append next into session
  emit TokenEvent(next)
  if stop_condition(next):
    emit FinalResponse(stop_token)
    break
\end{lstlisting}

真实服务还要处理取消和背压。客户端断开时，runtime 必须停止后续 decode，并释放该 session 的 KV cache、arena 和 pending events。若下游消费 token 很慢，不能让输出队列无限增长；要设置队列上限，超过后暂停调度、丢弃请求或返回 backpressure 错误。否则模型算得再快，服务也会被 I/O 队列拖垮。

\begin{warningbox}
流式输出是运行时协议，不是把每个 token 打印出来。它必须和 session 生命周期、KV cache 释放、取消、超时、trace 和 backpressure 绑定。
\end{warningbox}

\subsection{容量预算和 admission control}

业务服务最先撞到的不是“能不能生成一句话”，而是能同时服务多少请求。admission control 的职责是在请求进入 decode 前判断资源是否足够。最小预算包括权重、packed weight、KV cache、activation arena、workspace、队列和系统余量。

\begin{lstlisting}[numbers=none,caption={请求进入前的容量判断}]
required_kv_bytes =
  layers * max_context_tokens * kv_heads * head_dim *
  2 /* K and V */ * kv_bytes_per_element

required_runtime_bytes =
  required_kv_bytes +
  activation_arena_bytes +
  workspace_bytes +
  trace_buffer_bytes
\end{lstlisting}

若使用前文 7B 级 GQA 参数，\code{L=32}、\code{kv\_heads=8}、\code{head\_dim=128}、fp16 KV、\code{4096} context，一个 session 的 KV cache 约 \code{512 MiB}。若 int4 权重约 \code{3.7 GiB}，机器剩余可用内存 \code{10 GiB}，理论上不能简单做 \code{10 GiB / 512 MiB}。还要扣掉 packed weight、arena、workspace、allocator 碎片、线程栈、页表和操作系统余量。

admission decision 应该结构化：

\begin{lstlisting}[numbers=none,caption={AdmissionDecision 字段}]
AdmissionDecision:
  accepted: true | false
  reason:
    ok | context_too_long | memory_budget_exceeded |
    queue_overloaded | backend_unavailable
  selected_plan_id
  reserved_kv_bytes
  reserved_workspace_bytes
  estimated_ttft_ms
  estimated_decode_tokens_per_s
\end{lstlisting}

这让业务层能给出明确错误：上下文太长、内存不够、队列满、GPU 后端不可用，而不是笼统返回“生成失败”。也让运维能知道是容量问题还是正确性问题。

\subsection{SLO 要分首 token 和后续 token}

推理服务的延迟目标不能只写一个平均响应时间。至少要拆成：

\begin{itemize}
  \item \term{TTFT}{time to first token}：从请求进入到第一个 token 可输出。
  \item decode p50/p95：每个后续 token 的延迟分布。
  \item prefill tokens/s：prompt 处理吞吐。
  \item decode tokens/s：生成阶段吞吐。
  \item queue wait：进入模型前等待了多久。
  \item memory peak：请求期间峰值内存或显存。
\end{itemize}

prompt 很长但输出很短时，TTFT 主要由 tokenizer、prefill、KV 写入和首个 sampler 决定；prompt 较短但输出很长时，decode p95、KV 读流和 sampler 更重要。若服务端合批，平均 tokens/s 可能上升，但单请求 p95 可能变差，因为请求要等待同批次其他 session。

\begin{lstlisting}[numbers=none,caption={业务 SLO 报告字段}]
ServiceSloReport:
  prompt_bucket
  max_context_bucket
  backend_plan_id
  quant_profile_id
  queue_wait_p50_p95
  ttft_p50_p95
  decode_step_p50_p95
  tokens_per_second
  rejection_rate
  memory_peak_bytes
  fallback_count
\end{lstlisting}

有了这些字段，才谈得上把第三册内容落到一个本地助手、RAG 服务、代码补全服务或批量摘要服务里。不同业务目标会选择不同 profile，但底层证据字段是同一套。

\subsection{排队模型：吞吐够不等于 p95 够}

Serving 最容易被平均值骗。一个模型离线测得 \code{20 tok/s}，不代表线上每个请求都能稳定低延迟。请求到达有波动，prompt 长度有波动，decode 长度有波动，prefill 和 decode 还会争同一组后端资源。只要利用率接近 1，排队等待就会迅速放大，p95/p99 会先崩。

先用最小模型描述单个服务台：

\[
\rho = \lambda \times E[S]
\]

其中 \(\lambda\) 是请求到达率，\(E[S]\) 是平均服务时间，\(\rho\) 是利用率。若 \(\rho \ge 1\)，长期看队列必然增长；若 \(\rho\) 虽小于 1 但接近 1，稍微有长 prompt 或长输出，尾延迟就会明显上升。Little's Law 给出另一个必须检查的关系：

\[
L = \lambda \times W
\]

\(L\) 是系统中平均请求数，\(W\) 是请求在系统里的平均停留时间。这个式子不能直接告诉 p99，但能抓住一个硬约束：如果到达率不变而平均停留时间上升，系统内在途请求数一定上升；在推理服务里，这通常意味着 KV cache、队列内存和取消恢复压力一起上升。

LLM serving 还要把一次请求拆成 prefill 和 decode：

\begin{lstlisting}[numbers=none,caption={请求服务时间的最小拆分}]
service_time_ms =
  queue_wait_ms +
  tokenizer_ms +
  prefill_ms(prompt_tokens) +
  first_sample_ms +
  sum(decode_step_ms(context_i) for each generated token) +
  stream_flush_ms
\end{lstlisting}

这里的 \code{decode\_step\_ms(context\_i)} 不是常数。前面 KV 账本已经说明，context 越长，每步至少要读更多历史 K/V。因此，同样生成 128 个 token，短 prompt 和长 prompt 的 decode 时间分布不同；同样 prompt，输出越长，后面的 token 也会越来越贵。

\begin{center}
\small
\begin{tabular}{p{0.25\linewidth}p{0.30\linewidth}p{0.31\linewidth}}
\toprule
现象 & 排队模型解释 & 工程动作 \\
\midrule
平均 tokens/s 高，p95 TTFT 高 & prefill 在队列里等待合批或后端空闲 & 分离 prefill/decode 队列，限制 batch 等待 \\
decode p99 随并发突增 & 利用率接近 1，长输出占住 decode slot & admission control，max new tokens，抢占或分片 \\
长上下文请求拖慢全部请求 & 单步 decode 服务时间随 \(S\) 上升 & context 分桶调度，KV 预算隔离 \\
取消后内存仍高 & in-flight 请求未及时释放 KV/page & cancellation barrier，session epoch，回收审计 \\
队列满但 GPU 利用率不高 & CPU tokenizer、sampler、拷贝或锁成为服务台 & 分阶段 trace，拆服务台利用率 \\
\bottomrule
\end{tabular}
\end{center}

\topic{continuous batching 的取舍。} Continuous batching 把不同请求的 decode step 组合到同一批执行，提高后端利用率。它的代价是调度器每步都要决定哪些 session 进入 batch，如何处理刚完成 prefill 的请求，如何处理取消、超时和不同 context 长度。它改善吞吐，但可能增加单请求等待，尤其是短请求被长请求反复同批拖住时。

\topic{paged KV 的 SLO 意义。} Paged KV 不只是 kernel layout。它让 session 的 KV cache 可以按 block 分配、释放和复用，降低变长请求造成的碎片；也让 prefix cache 和迁移更容易。但它把一个连续地址公式换成 block table 查找，profile 必须同时记录 block 命中、page 分布、释放延迟和 slot epoch。否则，paged KV 的正确性 bug 会表现成偶发 logits 漂移，性能 bug 会表现成 p99 decode 抖动。

\topic{admission control 不是简单限流。} 接收请求前要估算两类资源：静态资源是否够，例如 KV 最大容量、workspace 和队列长度；动态资源是否会把利用率推到危险区，例如预计 prefill 时间、预计 decode token 数、当前 in-flight decode step。一个实用的 admission report 至少要多写几行：

\begin{lstlisting}[numbers=none,caption={带排队模型的 admission report}]
AdmissionEstimate:
  arrival_rate_1m
  inflight_requests
  queue_depth
  estimated_prefill_ms
  estimated_decode_ms_p50
  estimated_decode_ms_p95
  estimated_kv_bytes
  projected_utilization
  projected_ttft_p95
  decision: accept | reject | shed | retry_after
\end{lstlisting}

若报告没有 \code{projected\_utilization} 和 \code{queue\_depth}，它只能解释“内存够不够”，不能解释“接了这个请求后 p95 会不会炸”。若没有 \code{estimated\_kv\_bytes}，它只能解释排队，不能解释为什么取消恢复和长上下文会造成内存压力。

\subsection{错误也要沿链路定位}

链路完整的好处是错误能定位。若生成结果不对，不要直接怀疑 sampler 或模型质量，而要按链路缩小范围：

\begin{itemize}
  \item tokenizer 输出是否和 reference 一致。
  \item manifest 中 embedding、lm head、Q/K/V、MLP 权重角色是否正确。
  \item shape/dtype/quant verifier 是否通过。
  \item packed weight 与 reference converter 是否一致。
  \item 单算子 oracle 是否通过。
  \item 层输出 oracle 从哪一层开始漂移。
  \item logits top-k 是否变化。
  \item sampler 固定 seed 后是否仍然分叉。
\end{itemize}

性能问题也一样。首 token 慢，不等于 decode 慢；decode 慢，不等于所有 kernel 都慢；GPU kernel 时间低，不等于端到端快。profiler 要沿链路拆：load、prepare、prefill、decode、sample、copy、sync、fallback。

\subsection{用 trace 把链路变成证据}

链路如果只存在于图里，仍然不能维护。每次运行都应该能输出一个轻量 trace，让 correctness 和 performance 共享同一套定位坐标。trace 不必一开始很复杂，但字段要稳定。

\begin{lstlisting}[numbers=none,caption={端到端推理 trace 的最小字段}]
RunTrace:
  run_id
  model_id
  manifest_hash
  backend_plan_id
  quant_profile_id
  prompt_tokens
  generated_tokens
  events:
    stage
    segment_id
    layer_id
    backend
    context_length
    start_ns
    end_ns
    bytes_read_estimate
    bytes_written_estimate
    oracle_status
\end{lstlisting}

\code{manifest_hash} 追溯模型语义；\code{backend_plan_id} 追溯 kernel、layout 和 fallback；\code{quant_profile_id} 追溯误差阈值和量化格式；\code{context_length} 让 decode 指标可分档。trace 坐标稳定，后续 CPU SIMD、GPU kernel 和 KV cache 量化报告才可比较。

\subsection{硬验收}

读完本节，应该能用一条链路解释全书：

\begin{itemize}
  \item 模型资产先变成 manifest 和 typed graph，再 lowering 成 executable plan。
  \item 请求运行时绑定 token、KV cache、arena、workspace、sampler、oracle 和 profiler。
  \item prefill 与 decode 是两种形态，不能强行共用一个性能模型。
  \item 模型、编译、请求和算子生命周期必须分开，否则内存、状态和测量都会串扰。
  \item executable plan 要显式记录 segment、backend、workspace、状态读写、oracle checkpoint 和 profiler label。
  \item 正确性和性能问题都应该沿链路定位，而不是从最终文本倒猜。
  \item 端到端 trace 要让模型语义、后端计划、量化 profile、context length 和事件耗时可追溯。
\end{itemize}
